\documentclass[11pt]{article}
\usepackage[letterpaper,margin=0.95in]{geometry}
\usepackage[T1]{fontenc}
\usepackage{lmodern,amsmath,amssymb,amsthm,booktabs,microtype}
\usepackage[colorlinks=true,linkcolor=blue,citecolor=blue,urlcolor=blue]{hyperref}
\newtheorem{theorem}{Theorem}
\newtheorem{lemma}{Lemma}
\newtheorem{remark}{Remark}
\newcommand{\A}{\mathcal A}
\newcommand{\ii}{\mathrm i}
\newcommand{\PP}{\mathbb P}
\DeclareMathOperator{\sn}{sn}
\DeclareMathOperator{\cn}{cn}
\DeclareMathOperator{\dn}{dn}
\title{Verification of the odd-period\protect\\ focal inversion area product}
\author{Alec Kriebel\\[3pt]\small Independent researcher\\
\small\href{https://orcid.org/0009-0001-9320-500X}{ORCID: 0009-0001-9320-500X}}
\date{16 October 2026}
\hypersetup{pdftitle={Verification of the odd-period focal inversion area product},pdfauthor={Alec Kriebel}}
\begin{document}
\maketitle
\begin{abstract}
For an odd periodic billiard in a proper confocal ellipse pair, invert its
original vertices in unit circles centered at the original outer foci.
The product of the two signed straight-edge areas is independent of phase.
This is the historical invariant k903,a of Reznik--Garcia--Koiller and
Conjecture 1 of Reznik--Garcia--Helman. Ferudun's 2026 preprint proves the
same theorem using substantially the same meromorphic divisor mechanism.
This attributed verification note gives a self-contained intrinsic
flag-curve presentation: four isotropic common tangents identify every
possible pole, odd torsion separates their orbit grids, and reversal
supplies the opposite-focus zeros needed for compactness.
Primitive admitted star windings, signs, and indexed repetitions are
specified. Four exact finite billiard controls accompany the proof,
including two same-caustic even orbits with different products.
No novelty or absolute-priority claim is made.
\end{abstract}

\section{Historical target and attribution}

Reznik, Garcia, and Koiller define signed polygon area in
\cite[Eq.~(1)]{RGK}, original-vertex inversion about the original billiard
foci in \cite[\S3.9]{RGK}, and the odd focal-inverse area product in
\cite[Table~10, p.~12, k903,a]{RGK}. The row is dated 10/20. The shorter
journal edition does not contain this row. Reznik, Garcia, and Helman give
the triangle product as Proposition~7 and its odd-period extension as
Conjecture~1 in \cite[p.~6]{RGH}.

Ferudun's \emph{Products of the Two Focal Inversion Areas in Elliptic
Billiards}, version 1.0 \cite{Ferudun}, states and proves this same
strict odd-period theorem, including admitted stars, reversal, and
repetitions. Its Jacobi-coordinate proof and the flag-curve proof below
share the essential architecture: bound the area poles, use reversal to
reduce their orders, obtain opposite-focus zeros, and apply compactness.
The intrinsic formulation here is an attributed verification and
exposition of that theorem and mechanism.

Complexification, flag involutions, and removal of poles on a compact
curve are established tools for Poncelet invariants; see
Akopyan--Schwartz--Tabachnikov \cite{AST}, Ch\'avez-Caliz \cite{Chavez},
and Schwartz \cite{Schwartz}. Their cited invariants concern other
observables. For the present inverse areas the focal divisor and
adjacent-pole calculation must still be supplied; these are made explicit
below. Source and chronology limitations are recorded in Section~\ref{sec:limits}.

\section{Statement and conventions}

Fix
\begin{equation}\label{eq:pair}
 E:\frac{x^2}{a^2}+\frac{y^2}{b^2}=1,
 \qquad C:\frac{x^2}{\alpha^2}+\frac{y^2}{\beta^2}=1,
 \quad
 \begin{gathered}
 a>b>0,\quad 0<\lambda<b^2,\\
 \alpha^2=a^2-\lambda,\quad\beta^2=b^2-\lambda.
 \end{gathered}
\end{equation}
Put \(c=\sqrt{a^2-b^2}>0\), so \(\alpha^2-\beta^2=c^2\), and let
\(F_f=(f,0)\), \(f=\pm c\). For an ordered real billiard
\((P_0,\ldots,P_{N-1})\), with indices taken modulo \(N\), define
\begin{equation}\label{eq:inversearea}
 I_f(P)=F_f+\frac{P-F_f}{\|P-F_f\|^2},\qquad
 A_f=\frac12\sum_{j=0}^{N-1}\det(I_f(P_j),I_f(P_{j+1})).
\end{equation}
Only the vertices are inverted; successive image vertices are joined by
straight segments. For \(P=(a\cos\theta,b\sin\theta)\),
\(\|P-F_f\|=a-f\cos\theta\ge a-c>0\), so real inversion is nonsingular.

\begin{theorem}\label{thm:main}
Suppose the pair \eqref{eq:pair} admits a primitive real billiard of odd
period \(N\ge3\). For each oriented family of such billiards, the product
\(A_cA_{-c}\) in \eqref{eq:inversearea} is independent of phase and is
strictly positive. The assertion includes each coprime star winding
actually admitted by the pair. Reversal negates both areas and preserves
the product. Repeating a primitive cycle \(r\ge1\) times multiplies each
area by \(r\) and the product by \(r^2\).
\end{theorem}

Primitive means that the ordered billiard flag first returns after
\(N\) steps. The star formulation does not permit arbitrary reordering
of one vertex set. In Stachel's convention \cite[Theorem~4.3]{Stachel},
the modulus is \(k=c/\alpha\in(0,1)\), the vertices are
\((-a\sn u,b\cn u)\), and the step is \(2v\), where
\[
 \cn v=\beta/b,\qquad a=\alpha\dn v/\cn v,\qquad
 v=2\nu K(k)/N,\quad 1\le\nu<N/2,\quad\gcd(\nu,N)=1.
\]
The software parameter for these functions is \(k^2\). These relations
specify the period and winding conventions; the proof uses no complex
Jacobi period identity. Circular, coincident, degenerate, and hyperbolic
caustics are outside the strict statement.

\section{The compact flag curve}

Let \(X\) be the normalized projective curve of flags \((P,\ell)\), where
\(P\in E\) and \(\ell\) is a tangent to \(C\) through \(P\). Projection
\(\pi:X\to E\) has degree two. Its branch vertices are \(E\cap C\):
\begin{equation}\label{eq:branches}
 x^2=\frac{a^2\alpha^2}{c^2},\qquad
 y^2=-\frac{b^2\beta^2}{c^2}.
\end{equation}
They are four distinct finite points with nonzero coordinates. The conic
gradients there are independent, since
\(a^2/\alpha^2\ne b^2/\beta^2\).

An explicit check on normalization and infinity is useful. Parametrize
\(E\) by
\[
 x=a\frac{1-t^2}{1+t^2},\qquad y=\frac{2bt}{1+t^2}.
\]
The tangent discriminant gives the double cover, up to a nonzero square
factor, as
\begin{equation}\label{eq:quartic}
 q^2=\beta^2t^4+(4\alpha^2-2\beta^2)t^2+\beta^2.
\end{equation}
The discriminant in \(t^2\) is \(16\alpha^2c^2\ne0\), and the constant
and leading coefficients are nonzero. Thus the four branch roots are
distinct and nonzero. At \(t=\pm\ii\), the right side is
\(-4c^2\ne0\); these are the geometric points at infinity of \(E\).
The parameter value \(t=\infty\) instead lies over \((-a,0)\) and is
unbranched. The compact cover is connected, because it has simple branch
points, and has genus one by Riemann--Hurwitz.

Let \(\sigma\) exchange the two outer endpoints on \(\ell\), and let
\(\tau\) exchange the two caustic tangents through \(P\). These extend
to involutions of \(X\), with
\begin{equation}\label{eq:involutions}
 T=\tau\sigma,\quad \pi\tau=\pi,\quad \pi T=\pi\sigma,
 \qquad \sigma T\sigma=T^{-1}.
\end{equation}
The map \(T\) advances along the tangent chord and then chooses the
other outgoing caustic tangent: it is the billiard successor.
Each involution has fixed points and is nonidentity. A fixed-point
involution on a genus-one curve acts in an Abel coordinate as
\(u\mapsto d-u\): its linear part is \(-1\), since a nonidentity
translation cannot fix a point. Consequently \(T\) is a translation.
A primitive period-\(N\) real flag makes its translation element have
exact order \(N\) on all of \(X\).

A fixed flag of \(\sigma\) corresponds to a common tangent of the two
conics. For a line \(ux+vy+w_0=0\), the two tangency conditions are
\[
 a^2u^2+b^2v^2=w_0^2,\qquad
 \alpha^2u^2+\beta^2v^2=w_0^2.
\]
Their difference is \(\lambda(u^2+v^2)=0\). Thus the four common
tangents, all distinct and transverse as dual-conic intersections, are
\begin{equation}\label{eq:isotropic}
 z=x+\ii y=\pm c,\qquad w=x-\ii y=\pm c.
\end{equation}
Write their \(\sigma\)-fixed flags as \(s_{f,z},s_{f,w}\), \(f=\pm c\).

\section{Complete focal pole inventory}

Translation of every image vertex by \(-F_f\) leaves its closed
shoelace area unchanged. In the two complex coordinates, inversion then
has the meromorphic form
\begin{equation}\label{eq:reciprocal}
 Z_f=\frac1{w-f},\qquad W_f=\frac1{z-f}.
\end{equation}
These are algebraic extensions of ordinary real inversion; complex
conjugation is not being used as a holomorphic operation.
On \(E\),
\begin{equation}\label{eq:factor}
 z-f=-\frac{a+f}{1+t^2}
       \left(t-\frac{\ii b}{a+f}\right)^2.
\end{equation}
Hence \(W_f\) has a double pole at the isotropic contact
\[
 P_{f,z}=\left(\frac{a^2}{f},\frac{\ii b^2}{f}\right),
 \qquad w-f=\frac{2b^2}{f}\ne0.
\]
At \(P_{f,w}=(a^2/f,-\ii b^2/f)\), the roles of \(Z_f,W_f\) are
reversed. Each contact has caustic residual
\[
 \frac{x^2}{\alpha^2}+\frac{y^2}{\beta^2}-1
 =-\frac{\lambda^2}{\alpha^2\beta^2}\ne0.
\]
Therefore \(\pi\) is unramified there. Above \(P_{f,z}\) the two flags
are \(s_{f,z}\) and \(\tau s_{f,z}=Ts_{f,z}\); each pole has order two.
The analogous statement holds above \(P_{f,w}\).

There are no hidden poles at infinity. At either geometric infinity,
\(y=\pm\ii(b/a)x\), so both \(z\) and \(w\) have simple poles with
nonzero leading coefficients proportional to \(a\pm b\).
Their reciprocals in \eqref{eq:reciprocal} vanish. We have thus listed
every polar flag of both inverse coordinates.

\begin{lemma}\label{lem:grids}
If \(N\) is odd, the four grids
\(G_s=\{T^js:j\in\mathbb Z/N\}\), for the four fixed flags in
\eqref{eq:isotropic}, are pairwise disjoint.
\end{lemma}
\begin{proof}
If \(s'=T^js\) and both flags are fixed by \(\sigma\), then
\eqref{eq:involutions} gives \(s'=T^{-j}s\). Hence \(T^{2j}s=s\).
Exact order \(N\) and oddness give \(j=0\pmod N\), so \(s'=s\).
\end{proof}

Every possible area pole lies on one of these grids, because its
coordinate poles occur at \(s,Ts\). The proof of the lemma uses no
primality or generic-parameter assumption. It covers composite odd
periods and every admitted primitive winding.

\section{Reversal removes every product pole}

For a phase \(Q\in X\), put \(R_j=W_f(T^jQ)\), \(S_j=Z_f(T^jQ)\).
The complexified signed area is
\begin{equation}\label{eq:trace}
 \A_f(Q)=\frac1{4\ii}\sum_j(R_jS_{j+1}-S_jR_{j+1})
       =\frac1{4\ii}\sum_jR_j(S_{j+1}-S_{j-1}).
\end{equation}
It is invariant under \(T\). Moreover
\(\pi T^j\sigma Q=\pi T^{1-j}Q\), so \(\sigma\) reverses the vertex
list and
\begin{equation}\label{eq:odd}
 \A_f(\sigma Q)=-\A_f(Q).
\end{equation}

At \(s=s_{f,z}\), only vertices 0 and 1 have poles, both of order two
in the \emph{same} coordinate \(R\). Every \(S_j\) is regular there by
Lemma~\ref{lem:grids}. In particular, the determinant for the edge joining
the two singular vertices contains no \(R_0R_1\) term. The area pole
therefore has order at most two before any symmetry is applied.

Choose a local coordinate \(h\) at the fixed point \(s\) with
\(\sigma:h\mapsto-h\). Equation~\eqref{eq:odd} eliminates the double
coefficient and constant term, giving \(\A_f=\rho/h+O(h)\).
For an explicit principal-part check, write
\[
 R_0=r_{-2}h^{-2}+r_{-1}h^{-1}+O(1),\qquad
 D_0=S_1-S_{-1}=d_0+d_1h+O(h^2).
\]
Reversal gives \(R_1(h)=R_0(-h)\) and
\(D_1(h)=S_2-S_0=-D_0(-h)\). Thus the polar terms in \eqref{eq:trace}
are \([R_0(h)D_0(h)-R_0(-h)D_0(-h)]/(4\ii)\), and
\begin{equation}\label{eq:residue}
 \rho=\frac{r_{-2}d_1+r_{-1}d_0}{2\ii}.
\end{equation}
No unsupported trace-cancellation identity is needed.

At the same phase the opposite-focus area \(\A_{-f}\) is regular,
since its pole grids are disjoint from this one. It is odd by
\eqref{eq:odd}, so \(\A_{-f}=b_1h+O(h^3)\). The product is removable
and takes value \(\rho b_1\). At \(s_{f,w}\) interchange \(R,S\);
the same pole bound and oddness apply. Translation handles all points
of each grid. Elsewhere the inverse coordinates are regular. The
product is therefore holomorphic on compact connected \(X\), and is
constant. This proves phase constancy in Theorem~\ref{thm:main}.

Orient real tangent chords with \(C\) on their left. Each focus is
strictly inside \(C\), since \(\alpha^2-c^2=\beta^2>0\). Consequently
\[
 \det(P_j-F_f,P_{j+1}-F_f)>0.
\]
The corresponding determinant after inversion is this positive number
divided by
\(\|P_j-F_f\|^2\|P_{j+1}-F_f\|^2\). Both signed areas, and hence
their constant product \(K_0\), are positive. This argument is edgewise
and applies to stars. Reversal negates both shoelace sums. Repetition
simply repeats the sum \(r\) times, proving the stated \(r,r^2\) laws.

Positivity also sharpens the divisor description without circularity.
If \(\rho=0\) at any listed grid, both factors would be analytic and
odd there and their product would be zero. Thus all these poles and
opposite zeros are exactly simple; there are no additional zeros.
As reduced sums of \(N\) points per grid,
\[
 (\A_f)=G_{s_{-f,z}}+G_{s_{-f,w}}-G_{s_{f,z}}-G_{s_{f,w}}.
\]

\section{Exact finite controls and the parity restriction}

Take \(a^2=8\), \(b^2=5\), so \(f=\pm\sqrt3\). For
\(\lambda=40/9\), two positive primitive triangles are
\begin{align*}
 H&=\bigl((2\sqrt2,0),(-4\sqrt2/3,5/3),(-4\sqrt2/3,-5/3)\bigr),\\
 V&=\bigl((0,\sqrt5),(-8/3,-\sqrt5/3),(8/3,-\sqrt5/3)\bigr).
\end{align*}
Direct inversion and shoelace sums give
\[
 \begin{array}{c|cc|c}
 &A_c&A_{-c}&A_cA_{-c}\\ \hline
 H&\sqrt2/10-\sqrt3/80&\sqrt2/10+\sqrt3/80&5/256\\
 V&\sqrt5/16&\sqrt5/16&5/256
 \end{array}
\]
The source formula comparison requires care. On p.~3 of \cite{RGH},
the printed \(\delta=\sqrt{a^4-a^2b^2-b^4}\) has radicand \(-1\) at
these axes. Replacing its final minus by plus gives \(\delta=7\) and
makes the Proposition~7 expression agree with \(5/256\) at this exact
control. This checks that control only; no general corrected formula or
inspected journal correction is asserted. The universal proof above
uses neither expression.

Oddness is essential. For arbitrary \(a>b>0\), put
\(S=a^2+b^2\), \(\lambda=a^2b^2/S\), so
\(\alpha=a^2/\sqrt S\), \(\beta=b^2/\sqrt S\). Two genuine primitive
four-period billiards for this same pair are
\[
 D=((a,0),(0,b),(-a,0),(0,-b)),\qquad
 R=((\alpha,\beta),(-\alpha,\beta),(-\alpha,-\beta),(\alpha,-\beta)).
\]
For \(D\), its side lines \(\pm x/a\pm y/b=1\) are tangent to \(C\)
because \(\alpha^2/a^2+\beta^2/b^2=1\); reflection holds by axis
symmetry. For \(R\), the sides \(x=\pm\alpha,y=\pm\beta\) are
tangent, its vertices lie on \(E\), and the outer normal is proportional
to \((\pm1,\pm1)\), which reflects its horizontal and vertical edges.
Four distinct vertices establish primitivity.

For either focus, direct inversion gives \(A_f(D)=2/(ab)\) and
\(A_f(R)=1/(\alpha\beta)\). Their products differ:
\[
 A_c(D)A_{-c}(D)=\frac4{a^2b^2},\qquad
 A_c(R)A_{-c}(R)=\frac{S^2}{a^4b^4},\qquad
 \frac{S^2-4a^2b^2}{a^4b^4}=\frac{c^4}{a^4b^4}>0.
\]
At the specified axes and \(\lambda=40/13\), the rectangle vertices
are \((\pm8/\sqrt{13},\pm5/\sqrt{13})\) in cyclic order. The two
products are \(1/10\) and \(169/1600\). These even controls were
derived during the present verification and carry no discovery claim.

The accompanying standard-library exact verifier checks the four
polygons' incidences, segment tangencies, reflections, inversions,
areas, reversal, and triple repetition. It also distinguishes inversion
radius scaling from physical similarity: radius 2 multiplies each
area by 16 and the product by 256; scaling the orbit and both foci by
\(s=3/2\), with radius fixed at 1, multiplies area by \(s^{-2}\) and
product by \(s^{-4}\). Its 245 checks in 27 nonempty groups are finite
controls, not a certificate of the all-period proof. Explicit exceptions
remain active under Python optimization; deliberate false and malformed
inputs are tested separately. All four finite controls are convex
winding-one trajectories; higher stars are covered by the analytic
proof, not by this finite suite.

\section{Scope, record dates, and assistance}\label{sec:limits}

The priority comparison was bounded. In particular the official
15-page preview of \cite{Book} lists \S8.5, ``Invariant area product,''
but omits that section's text; the full section was unavailable for this
comparison. The RGH journal full text and its typography were not
inspected. These access gaps do not establish absence of earlier proofs.
The direct exact theorem and overlapping mechanism in \cite{Ferudun}
are expressly credited.

The preserved \href{https://github.com/AlecKriebel/Math/pull/206}{GitHub
PR206} record reports creation at 2026-10-01T05:36:26Z; its preserved
HEAD has committer time 2026-10-01T05:36:15Z. Ferudun's Zenodo record
reports creation at 2026-10-02T07:08:59.065762+00:00, and the paper
declares 2 October 2026 as its date. These are facts about those records.
They do not certify first public availability of every byte, visibility
history, intellectual independence, literal simultaneity, or absolute
priority. The note makes no claim that the historical conjecture was
still open when this verification was prepared.

AI tools were used extensively in research, proof development,
drafting, computation, and verification. Separate AI adversarial checks
are not independent human peer review. This is an unrefereed research
note; no conventional human refereeing or proof-assistant certification
is claimed.

\paragraph{License and materials.}
Copyright 2026 Alec Kriebel. The authored paper and accompanying
verification materials are available under
\href{https://creativecommons.org/licenses/by/4.0/}{CC BY 4.0}.
External works retain their own rights. The package includes source,
finite exact code, raw run receipts, and source/attribution notes;
it includes no third-party full texts or software dependencies.

\begin{thebibliography}{9}
\bibitem{RGK} D. Reznik, R. Garcia, J. Koiller,
\emph{Eighty New Invariants of N-Periodics in the Elliptic Billiard},
arXiv:2004.12497v11 (29 October 2020),
\url{https://arxiv.org/abs/2004.12497v11}.
\bibitem{RGH} D. Reznik, R. Garcia, M. Helman,
\emph{The Talented Mr. Inversive Triangle in the Elliptic Billiard},
arXiv:2012.03020v2 (13 July 2021),
\url{https://arxiv.org/abs/2012.03020v2}.
\bibitem{Ferudun} A. Ferudun,
\emph{Products of the Two Focal Inversion Areas in Elliptic Billiards},
version 1.0 (2 October 2026),
\url{https://doi.org/10.5281/zenodo.23096994}.
\bibitem{AST} A. Akopyan, R. Schwartz, S. Tabachnikov,
\emph{Billiards in ellipses revisited}, arXiv:2001.02934v2 (2020),
\url{https://arxiv.org/abs/2001.02934v2}.
\bibitem{Chavez} A. C. Ch\'avez-Caliz,
\emph{More About Areas and Centers of Poncelet Polygons},
Arnold Mathematical Journal \textbf{7} (2021), 91--106,
\url{https://doi.org/10.1007/s40598-020-00154-8}.
\bibitem{Schwartz} R. Schwartz,
\emph{Monodromy and Poncelet Families} (9 August 2008),
\url{https://www.math.brown.edu/reschwar/Papers/poncelet.pdf}.
\bibitem{Stachel} H. Stachel,
\emph{On the motion of billiards in ellipses},
European Journal of Mathematics \textbf{8} (2022), 1602--1622,
\url{https://doi.org/10.1007/s40879-021-00524-2}.
\bibitem{Book} R. Garcia, D. Reznik,
\emph{Poncelet Surprises in the Euclidean Plane}, IMPA (2021),
ISBN 978-65-89124-43-6; official preview,
\url{https://coloquio33.impa.br/pdf/33CBM06-eBook-preview.pdf}.
\end{thebibliography}
\end{document}
